\clearpage
\par\bigskip
\pdfbookmark[1]{Bemerkungen}{purae-bemerkungen}
\begin{center}\textsc{Bemerkungen.}\end{center}
\begingroup\small
Aus der ursprünglichen Fortsetzung dieser Abhandlung von Artikel~19 an, die der Behandlung besonderer Fälle gewidmet war, sind nur noch einige Artikel erhalten. Sie befassen sich mit der quadratischen Gleichung, deren Wurzeln die beiden Perioden mit $\frac{n-1}2$ Gliedern sind; das Manuskript bricht zu Beginn der Untersuchung ab, durch die das Vorzeichen der bei ihrer Lösung auftretenden Quadratwurzel bestimmt werden sollte. Die Übereinstimmung dieses noch erhaltenen Anfangs mit der Abhandlung Summatio quarumdam serierum singularium zeigt, dass der Verfasser seinen Plan änderte, um die eben genannte Bestimmung des Vorzeichens zum Gegenstand einer besonderen Abhandlung zu machen. Vergleicht man damit das Zitat in Artikel~8, wo im Manuskript versehentlich die dritte statt der zweiten Ausgabe des Werkes von \textsc{Lagrange} angegeben war, so ergibt sich, dass diese Handschrift aus dem Jahr 1808 stammt. Dass die Veröffentlichung des Vorangehenden nicht aufgegeben worden war, zeigt ein bei Artikel~19 offenbar später eingefügtes Blatt, auf dem eine andere Fortsetzung beginnt und hinsichtlich der Bestimmung des Vorzeichens bereits auf die Abhandlung Summatio usw. verwiesen wird. Diese zweite Fortsetzung, die aber ebenfalls bald abbricht, ist hier wiedergegeben. Der Text des durchweg druckfertigen Manuskripts wurde bei der Herausgabe getreu beibehalten; nur in Artikel~16 mussten die Formeln für den zweiten Fall hinzugefügt werden.
\begin{flushright}R. \textsc{Dedekind}.\end{flushright}
\endgroup
\BookRule{.34\textwidth}
